\setlength{\baselineskip}{16pt}
\chapter{作者简历及攻读博士学位期间取得的研究成果}\zihao{5}
\setlength{\parindent}{0pt}

%[内容采用五号宋体]  包括教育经历、工作经历、攻读学位期间发表的论文和完成的工作等。行距16磅，段前后各为0磅。
% \vspace{10pt}

郝赫，男，汉族，中共党员，1999年11月生，河北沧州人。

\vspace{10pt}
\textbf{一、教育经历}
\vspace{5pt}

2022年9月至今，北京交通大学，交通运输规划与管理专业，博士研究生（直接攻博）；

2025年9月至2026年4月，新加坡国立大学，国家公派联合培养博士研究生；

2018年9月至2022年6月，北京交通大学，交通工程专业，工学学士学位，获免试推荐直接攻读博士学位资格。

\vspace{10pt}
\textbf{二、期刊论文}
\vspace{5pt}

发表/在投期刊论文21篇（SCI/SSCI期刊14篇，EI期刊3篇，中文核心及其他期刊4篇），其中第一/通讯作者14篇，获ESI高被引论文热点论文1篇。

\begin{enumerate}[leftmargin=*, parsep=5pt, label={[\arabic*]}]

	\item \textbf{He Hao}, Enjian Yao*, Long Pan, Rongsheng Chen, Yue Wang, Hui Xiao. (2025). Exploring heterogeneous drivers and barriers in MaaS bundle subscriptions based on the willingness to shift to MaaS in one-trip scenarios[J]. \textbf{Transportation Research Part A: Policy and Practice}, 199: 104525. (\textbf{中科院1区Top，ESI高被引论文热点论文})

	\item Enjian Yao, \textbf{He Hao}*, Long Pan*, Rongsheng Chen, Yue Wang, Hui Xiao. (2025). Investigating the willingness of shifting to MaaS in one-trip scenarios: Insights from comparative stated surveys[J]. \textbf{Transportation Research Part A: Policy and Practice}, 192: 104384. (\textbf{中科院1区Top})

	\item \textbf{He Hao}, Enjian Yao*, Long Pan, Yang Yang, Yue Wang. (2026). Understanding travel behavior change with megacity development: Insights from decade-long household travel surveys in Beijing[J]. \textbf{Journal of Transport Geography}, 131: 104565. (\textbf{中科院1区Top})

	\item \textbf{He Hao}, Enjian Yao, Rongsheng Chen*, Yang Yang, Yue Wang, Wan Liu. (2026). AR2G-AIRL: A context-adaptive game-inverse reinforcement learning algorithm for modeling pedestrian-vehicle interaction considering group behaviors[J]. \textbf{IEEE Transactions on Intelligent Transportation Systems}. (\textbf{中科院1区Top})

	\item \textbf{He Hao}, Enjian Yao*, Rongsheng Chen, Long Pan, Yue Wang. (2024). Dynamic passenger route guidance in the multimodal transit system with graph representation and attention based deep reinforcement learning[J]. \textbf{IEEE Transactions on Intelligent Transportation Systems}. DOI: 10.1109/TITS.2024.3390172. (\textbf{中科院1区Top})

	\item \textbf{He Hao}, Enjian Yao*, Yang Yang, Shasha Liu, Long Pan, Yue Wang. (2024). How to build vibrant communities by utilizing functional zones? A community-detection-based approach for revealing the association between land use and community vibrancy[J]. \textbf{Cities}, 155: 105431. (\textbf{中科院1区Top})

	\item \textbf{He Hao}, Enjian Yao*, Rongsheng Chen, Long Pan, Shasha Liu, Yue Wang, Hui Xiao. (2024). An approach for evaluating added values of MaaS bundles considering heterogeneous subscription willingness[J]. \textbf{Transportation}: 1--35. (\textbf{学院认定A+顶级期刊})

	\item 姚恩建, 王月, \textbf{郝赫}*, 杨扬, 郇宁. (2026). 城市群综合交通可达性对城际出行需求的影响研究[J]. \textbf{中国公路学报}, 39(5): 387--399. (\textbf{EI，学院认定A类权威期刊})

	\item 姚恩建, 陈卓利, \textbf{郝赫}*, 陈荣升, 杨扬. (2025). 基于强化学习的自动驾驶车辆路上突发障碍物换道避障控制算法[J]. \textbf{北京交通大学学报}, 49(5): 82--93. (\textbf{学院认定B类鼓励期刊})

	\item \textbf{郝赫}, 姚恩建*, 陈荣升, 潘龙, 王月. (2025). 基于深度学习的多模式公交线网乘客路径诱导策略研究[J]. \textbf{交通工程}, 25(2): 1--10. (\textbf{公路运输领域高质量科技期刊})
	
	\item \textbf{He Hao}, Enjian Yao*, Rongsheng Chen, Yue Wang, Yang Yang, Hui Xiao. (2026). From Single Trip Choice to Bundle Subscription: Understanding the Behavioral Transitions and Determinants in Two-stage MaaS Adoption Using a Jointly Estimated Model[J]. \textbf{Transportation Research Part A: Policy and Practice}. (\textbf{Major Revision，中科院1区Top})
	
	\item Enjian Yao, Yue Wang, \textbf{He Hao}*, Yang Yang*. (2026). Assessing the Alignment between Potential and Experienced Accessibility within Urban Agglomerations[J]. \textbf{Journal of Transport Geography}. (\textbf{Major Revision，中科院1区Top})
	
	\item \textbf{He Hao}, Qiang Meng, Enjian Yao*, Yang Yang. (2026). CoHACT: A Deep Generative Framework for Coordinated Household Activity-Chain Generation with Intra-Household Travel[J]. \textbf{Transportation Research Part B: Methodological}. (\textbf{Under Review，中科院1区Top})

	\item \textbf{He Hao}, Qiang Meng, Enjian Yao*, Yang Yang. (2026). Large-scale Household-level Population Synthesis: An Effective Hierarchical Deep Learning-based Model[J]. \textbf{Transportation Research Part B: Methodological}. (\textbf{Under Review，中科院1区Top})
	
	\item Ning Huan, Jiaoe Wang, Toshiyuki Yamamoto, \textbf{He Hao}, Hanqi Zhao, Enjian Yao*. (2026). Vehicle energy type preferences with hydrogen fuel cell vehicles as an alternative in electrified mobility. \textbf{Transportation Research Part D: Transport and Environment}, 159, 105500.

	\item Yue Wang, Enjian Yao*, Yongsheng Zhang, Long Pan, \textbf{He Hao}. (2024). Deep learning model for short-term origin-destination distribution prediction in urban rail transit network considering destination choice behavior[J]. \textbf{Transportation Research Record}: 03611981241243081.

	\item Wan Liu, Jinguang Liu, Shuai Dai, \textbf{He Hao}. (2026). A screening and prioritization method for urban road hazard points: Large-scale validation analysis in 18 cities[J]. \textbf{IEEE Access}.
	
	\item Long Pan, Enjian Yao, Yang Yang, \textbf{He Hao}. (2025). Planning of energy supply facilities for highway networks: Methodology research and practical applications[J]. \textbf{Transportation Research Procedia}, 92: 53.

	\item 王月, 姚恩建*, \textbf{郝赫}, 李义罡, 史佳柠. (2025). 京津冀城市群总体与城际出行特征异同性挖掘[J]. \textbf{地理学报}, 80(4): 1089--1102. (EI)

	\item 王月, 姚恩建*, \textbf{郝赫}. (2023). 低碳导向的多模式交通出行服务定价策略[J]. \textbf{清华大学学报（自然科学版）}. DOI: 10.16511/j.cnki.qhdxxb.2023.26.024. (EI)

	\item 王学勇, 姚恩建*, 王小华, \textbf{郝赫}. (2025). 我国城市停车配建标准体系研究[J]. \textbf{综合运输}, 47(2): 22--25.
	
\end{enumerate}

\vspace{10pt}
\textbf{三、会议论文}
\vspace{5pt}

发表/录用在投会议论文8篇，其中以第一作者或通讯作者身份完成3篇，相关成果在TRB、WTC、ITSWC及TRS等国际会议进行交流。
	
\begin{enumerate}[leftmargin=*, parsep=5pt, label={[\arabic*]}]

	\item \textbf{He Hao}, Enjian Yao*, Long Pan, Yang Yang, Rongsheng Chen, Yue Wang. (2025). Understanding travel behavior change and the role of public transport with megacity development[C]. \textbf{Transportation Research Board (TRB) 104th Annual Meeting}, Washington, D.C., USA.

	\item \textbf{He Hao}, Enjian Yao*, Long Pan, Yang Yang, Rongsheng Chen, Yue Wang. (2025). What does 2023 household travel survey in Beijing tell us? The roles of innovative survey methods and changes in travel mode preferences of residents[C]. \textbf{Transportation Research Board (TRB) 104th Annual Meeting}, Washington, D.C., USA.
	
	\item Rui Chen, Enjian Yao, \textbf{He Hao}*. (2026). Investigating the causal impact of subway expansion on carbon emissions in Beijing, China[C]. \textbf{World Conference on Transport Research (WCTR 2026)}, Toulouse, France.

	\item 王昊然, 姚恩建, 陈荣升, \textbf{郝赫}. (2026). 基于紧迫度公平的混合交通流自动驾驶协同变道策略[C]. \textbf{2026世界交通运输大会（WTC 2026）}. 

	\item 胡雪逸, 姚恩建, 姚辉, \textbf{郝赫}. (2026). 基于改进YOLOv8算法的沥青路面裂缝检测模型[C]. \textbf{2026世界交通运输大会（WTC 2026）}. 

	\item Long Pan, Enjian Yao, Yang Yang, \textbf{He Hao}. (2025). Planning of energy supply facilities for highway networks: Methodology research and practical applications[C]. \textbf{Transportation Research Symposium 2025}, Rotterdam, the Netherlands.

	\item Yue Wang, Enjian Yao, Yongsheng Zhang, Long Pan, \textbf{He Hao}. (2024). A deep learning model for short-term origin-destination distribution prediction in urban rail transit network considering destination choice behavior[C]. \textbf{Transportation Research Board (TRB) 103rd Annual Meeting}, Washington, D.C., USA.

	\item Yue Wang, Enjian Yao, Yongsheng Zhang, \textbf{He Hao}, Long Pan. (2023). Urban rail transit short-term OD flow prediction considering temporal-spatial characteristics and probability[C]. \textbf{29th ITS World Congress}, Suzhou, China.

\end{enumerate}

\vspace{10pt}
\textbf{四、发明专利}
\vspace{5pt}

获授权国家发明专利2件：

\begin{enumerate}[leftmargin=*, parsep=5pt, label={[\arabic*]}]

	\item 姚恩建, \textbf{郝赫}, 陈荣升, 杨扬, 潘龙, 王月. (2025). 基于群体特征的人车交互下行人意图判断与轨迹生成方法, 专利号：ZL202510378056.9, 公告号：CN119888673B.

	\item 姚恩建, \textbf{郝赫}, 潘龙, 杨扬, 王月. (2025). 基于居民出行影响因素与交通资源需求的预测分类方法, 专利号：ZL202510838732.6, 公告号：CN120409832B.

\end{enumerate}

\vspace{10pt}
\textbf{五、软件著作权}
\vspace{5pt}

登记软件著作权2件：

\begin{enumerate}[leftmargin=*, parsep=5pt, label={[\arabic*]}]

	\item 兰斌, 姚恩建, \textbf{郝赫}等. 故障大数据分析系统V1.0, 登记号：2025SR1812113.

	\item 兰斌, 姚恩建, \textbf{郝赫}等. 故障大数据挖掘系统V1.0, 登记号：2025SR1812140.

\end{enumerate}

\vspace{10pt}
\textbf{六、智库与著作}
\vspace{5pt}

撰写省部级智库建议2篇，参与编写国家重点出版物科普图书1部：

\begin{enumerate}[leftmargin=*, parsep=5pt, label={[\arabic*]}]

	\item 完善城市公交市场化运营机制. 姚恩建，\textbf{郝赫}，2024年被中央广播电视总台内参中心采用，并获中央广播电视总台感谢信1封。

	\item 加强多模式公交线网融合，推进北京市公交高质量发展. 姚恩建，孙迅，\textbf{郝赫}，2023年入选《交大智库建言》第42期，并获北京市交通委员会采纳。

	\item 姚恩建, 陈荣升等. (2026). 智慧的城市交通[M]. 北京: 人民交通出版社.（``十四五''国家重点出版物出版规划项目、中国科协``科普中国创作出版扶持计划''项目）\textbf{主要参编人员}

\end{enumerate}

\vspace{10pt}
\textbf{七、研究经历}
\vspace{5pt}

主持研究生创新项目1项，参与国家级、省部级及企事业单位委托科研项目17项，具体如下：

\begin{enumerate}[leftmargin=*, parsep=5pt, label={[\arabic*]}]

	\item 出行即服务理念下分级用户需求建模与运营优化理论研究. 中央高校基本科研业务费专项资金研究生创新项目（一类）, 2024.04--2026.03, 主持.

	\item 不同自主化水平道路交通系统风险辨识与安全管控技术. 国家重点研发计划“交通载运装备与智能交通技术”重点专项课题, 2023.12--2027.11, 参加.

	\item 大客流下综合客运枢纽多方式运力调控与韧性服务. 国家自然科学基金重点项目, 2025.01--2029.12, 参加.

	\item 城市群尺度下城际出行行为分析与交通供给优化. 国家自然科学基金面上项目, 2022.01--2025.12, 参加.

	\item 城轨运营突发事件下诱导信息对乘客行为影响及客流演化机理研究. 北京市自然科学基金轨道交通联合基金项目, 2022.01--2024.12, 参加.

	\item 面向民用机场的新能源设备规划与应用研究. 中国民用航空局民航安全能力建设基金, 2023.04--2025.12, 参加.

	\item 飞行区光伏可行性安全指标及相关政策机制研究. 中国民用航空局民航安全能力建设基金, 2024.07--2025.12, 参加.
	
	\item 京津冀跨界交通基础设施规建管养一体化机制研究. 首都高端智库决策咨询重大课题（省部级）, 2023.07--2025.12, 参加.

	\item 我市轨道交通与地面公交、慢行系统多网融合发展问题研究. 首都高端智库决策咨询重大课题（省部级）, 2022.03--2022.12, 参加.

	\item “四网融合”背景下的北京市郊铁路现代化运营管理体系研究. 首都高端智库决策咨询重点课题（省部级）, 2022.07--2023.03, 参加.

	\item 北京轨道交通与地面公交融合发展研究. 首都高端智库课题, 2025.11--2026.11, 参加.

	\item 北京市交通系统韧性提升研究. 首都高端智库课题, 2025.07--2026.03, 参加.

	\item 北京市郊铁路郊区站点周边交通及区域微循环优化研究. 首都高端智库课题, 2024.06--2026.08, 参加.

	\item 提高北京45分钟通勤比例可行措施研究. 首都高端智库决策咨询专项课题, 2022.11--2023.09, 参加.

	\item 需求导向下的“轨道+公交+慢行”三网融合规划方法研究. 交通运输部科学研究院, 2023.11--2025.08, 参加.

	\item 居民出行调查样本分析. 北京交通发展研究院, 2023.10--2027.09, 参加.

	\item 地铁故障大数据挖掘技术方法研究项目. 北京市地铁运营有限公司技术创新研究院分公司, 2023.09--2023.12, 参加.

	\item 北京市轨道交通13号线扩能提升工程道路交通工程方案咨询. 北京城建设计发展集团股份有限公司, 2023.02--2023.12, 参加.

\end{enumerate}

\vspace{10pt}
\textbf{八、荣誉奖励}
\vspace{5pt}

\makebox[7em][l]{2025年12月}中国科协青年科技人才培育工程博士生专项计划

\makebox[7em][l]{2026年5月}北京交通大学研究生知行奖学金（原``校长奖学金''）

\makebox[7em][l]{2025年12月}宝钢优秀学生奖学金（全校研究生仅3人）

\makebox[7em][l]{2025年12月}博士研究生国家奖学金

\makebox[7em][l]{2024年12月}博士研究生国家奖学金

\makebox[7em][l]{2025年5月}国家留学基金委联合培养博士生奖学金

\makebox[7em][l]{2026年6月}北京市优秀毕业生

\makebox[7em][l]{2026年3月}北京市三好学生

\makebox[7em][l]{2026年5月}关键软件技术创新成果（全国“十大”软件，示范性软件学院联盟）

\makebox[7em][l]{2024年12月}北京交通工程学会有奖征文一等奖（全国仅2项）

\makebox[7em][l]{2025年10月}第六届智慧交通大赛全国三等奖

\makebox[7em][l]{2024年11月}第二届“先导杯”智能车联网开放数据挑战赛全国二等奖

\makebox[7em][l]{2025年7月}世界交通运输大会（WTC）推荐论文

\makebox[7em][l]{2022年10月}“华为杯”第四届中国研究生人工智能创新大赛校二等奖

\makebox[7em][l]{2025年12月}北京交通大学三好学生


